% Звіт за етапом 4 — ШАБЛОН. Заповнити фактами (див. docs/stage4.md).
% Червоні [TODO] у PDF = незаповнені місця. Перед здачею їх не має бути.
\documentclass[12pt]{article}
% Спільна преамбула для звітів (XeLaTeX)
\usepackage{fontspec}
\usepackage{polyglossia}
\setmainlanguage{ukrainian}
\setotherlanguage{english}
\setmainfont{DejaVuSerif}[Extension=.ttf,UprightFont=*,BoldFont=*-Bold,ItalicFont=*-Italic,BoldItalicFont=*-BoldItalic]
\setsansfont{DejaVuSans}[Extension=.ttf,UprightFont=*,BoldFont=*-Bold,ItalicFont=*-Oblique,BoldItalicFont=*-BoldOblique]
\setmonofont{DejaVuSansMono}[Extension=.ttf,UprightFont=*,BoldFont=*-Bold,ItalicFont=*-Oblique,BoldItalicFont=*-BoldOblique,Scale=0.9]
\usepackage[a4paper,margin=2.2cm]{geometry}
\usepackage{booktabs}
\usepackage{tabularx}
\usepackage{array}
\usepackage{enumitem}
\usepackage{xcolor}
\usepackage{tcolorbox}
\usepackage{graphicx}
\usepackage{setspace}
\usepackage[hidelinks]{hyperref}
\usepackage{titlesec}
\usepackage{fancyhdr}
\usepackage{multirow}
\usepackage{colortbl}
\setlength{\parskip}{4pt}
\setlength{\parindent}{0pt}
\onehalfspacing

\definecolor{accent}{HTML}{1F4E79}
\definecolor{sgreen}{HTML}{E2F0D9}
\definecolor{wred}{HTML}{FBE5D6}
\definecolor{oblue}{HTML}{DEEBF7}
\definecolor{tyellow}{HTML}{FFF2CC}
\definecolor{lightgray}{HTML}{F2F2F2}

\titleformat{\section}{\large\bfseries\color{accent}}{\thesection}{0.6em}{}
\titleformat{\subsection}{\normalsize\bfseries}{\thesubsection}{0.6em}{}

\pagestyle{fancy}
\fancyhf{}
\lhead{\small ІТ, Модуль 2 — командний проєкт}
\rhead{\small Петров Б., Мекшаков А.}
\cfoot{\thepage}
\renewcommand{\headrulewidth}{0.3pt}

\newcolumntype{L}[1]{>{\raggedright\arraybackslash}p{#1}}
\newcolumntype{Y}{>{\raggedright\arraybackslash}X}

% Титульна сторінка
\newcommand{\maketitlepage}[3]{%
\begin{titlepage}
\centering
{\large Київський національний університет імені Тараса Шевченка\\[2pt]
Факультет комп'ютерних наук та кібернетики\par}
\vspace{3cm}
{\large Дисципліна: \textbf{Інформаційні технології}\par}
\vspace{0.5cm}
{\large Модуль 2 — робота в команді\par}
\vspace{2cm}
{\LARGE\bfseries #1\par}
\vspace{0.6cm}
{\large #2\par}
\vfill
\begin{flushright}
\begin{tabular}{l}
\textbf{Виконали:} студенти 4 курсу, група МІ-41\\[4pt]
Петров Богдан\\
Мекшаков Артемій\\[10pt]
\textbf{Спеціальність:} Комп'ютерні науки\\
\end{tabular}
\end{flushright}
\vspace{1.5cm}
{Київ — #3\par}
\end{titlepage}
}

\usepackage{longtable}
\usepackage{pdflscape}
\usepackage{float}
\graphicspath{{img/}{../docs/uml/out/}}
\newcommand{\todo}[1]{\textcolor{red}{\textbf{[TODO: #1]}}}
% Вставка зображення або TODO, якщо файлу ще немає
\newcommand{\fig}[3][0.9\textwidth]{%
\begin{figure}[H]\centering
\IfFileExists{img/#2}{\includegraphics[width=#1,height=0.7\textheight,keepaspectratio]{#2}}{%
\IfFileExists{../docs/uml/out/#2}{\includegraphics[width=#1,height=0.7\textheight,keepaspectratio]{#2}}{%
\fbox{\parbox{0.8\textwidth}{\centering\vspace{1.5cm}\todo{#2}\vspace{1.5cm}}}}}
\caption{#3}\end{figure}}
% Вставка згенерованої таблиці або TODO
\newcommand{\gen}[1]{\IfFileExists{generated/#1}{\input{generated/#1}}{\todo{generated/#1 --- запустити pnpm test:report}}}

\begin{document}
\maketitlepage{Звіт за етапом 4}{Реалізація прототипу: UML і схема БД, архітектура, тестування, API-документація}{2026}

\section{Мета етапу}
Реалізувати прототип StudyFlow в обсязі спринтів 1--3 (етап 3), промоделювати систему в UML, обґрунтувати архітектуру та інструментарій, покрити код модульними та інтеграційними тестами, задокументувати API і провести рефлексію спринта.

\section{План спринта і його виконання}
\todo{таблиця план (етап 3) vs факт по ID беклогу: ID, назва, SP, статус (Done / перенесено у спринт 4); підсумок SP план/факт}

\fig{notion-sprints.png}{Дошка спринтів у Notion}

\section{Поточна версія системи}
\todo{2--3 речення: що вміє система зараз}
\fig{dashboard.png}{Дашборд: найближчі дедлайни}
\fig{calendar-week.png}{Календар тижня}
\fig{group.png}{Сторінка групи: групові дедлайни та прогрес}
\fig{board.png}{Канбан командних задач}
\fig{mailpit.png}{Лист-нагадування (Mailpit, середовище розробки)}

\section{Архітектура системи та інструментарій}
\subsection{Архітектура}
\todo{модульний моноліт; шари Controller -- Service -- Prisma; модулі; same-origin через rewrites Next.js}
\fig{component.png}{Діаграма компонентів}
\fig{deployment.png}{Діаграма розгортання}

\subsection{Обґрунтування інструментарію}
\todo{таблиця: інструмент --- роль --- чому --- відкинута альтернатива (Next.js, NestJS, PostgreSQL, Prisma, JWT, Tailwind/shadcn, TanStack Query, Jest, supertest, Playwright, Docker Compose, GitHub Actions, Mailpit, PlantUML, Notion)}

\section{UML-моделювання}
\fig{use-case.png}{Діаграма варіантів використання}
\fig{class.png}{Діаграма класів доменної моделі}
\fig{sequence-group-deadline.png}{Послідовність: груповий дедлайн}
\fig{sequence-reminders.png}{Послідовність: надсилання нагадувань}
\fig{state.png}{Діаграма станів задачі / дедлайну}

\section{Схема бази даних}
\fig{er.png}{ER-діаграма бази даних}
\todo{таблиця сутностей: таблиця --- призначення --- ключові поля --- зв'язки}

\section{Модульні тести}
\todo{підхід: Jest, що мокається; підсумок}
\gen{summary.tex}
\gen{coverage.tex}
Повний журнал --- Додаток А.

\section{Інтеграційні та E2E-тести}
\todo{підхід: тестова БД studyflow\_test, stub пошти, Playwright}
\gen{integration.tex}
\gen{e2e.tex}

\section{API-документація}
\todo{Swagger UI на /api/docs, openapi.json у репозиторії; кількість ендпоінтів і тегів}
\fig{swagger.png}{Swagger UI}
Повний перелік ендпоінтів --- Додаток Б.

\section{Рефлексія спринта}
\todo{лише факти (git log, Notion, тести). Успіхи / Неуспіхи / Причини / Висновки і дії на спринт 4}

\section{Висновки і наступні кроки}
\todo{підсумок + етап 5: деплой (T-09), US-24, US-28, перенесене}

\appendix
\section{Журнал модульних тестів}
\gen{unit.tex}

\section{Перелік ендпоінтів API}
\todo{таблиця: метод, шлях, доступ, опис, ID беклогу}

\section{Інструкція запуску}
\todo{docker compose up -d; pnpm install; pnpm db:reset; pnpm dev; демо-акаунти}

\end{document}
